\documentclass[fleqn]{ctexart}
\usepackage{graphicx}
\usepackage{xcolor}
% 数学支持（提供 \mathbb、\dfrac 等）
\usepackage{amsmath,amssymb}
% 页面与版式设置：更小的页边距与顶部空白
\usepackage[a4paper,top=15mm,bottom=20mm,left=12mm,right=12mm]{geometry}
% 自动换行与字偶距优化
\usepackage{microtype}
\microtypesetup{tracking=true,protrusion=true,final}
% 在需要时放宽断行以避免溢出（数值可按需调整）
\setlength{\emergencystretch}{2em}
% 列表控制，便于让(1)后直接跟内容
\usepackage{enumitem}
% verbatim 环境支持（用于直接粘贴源码与输出）
\usepackage{verbatim}
\usepackage{fancyvrb}
\usepackage{float}
\usepackage{fvextra}
\RecustomVerbatimEnvironment{verbatim}{Verbatim}{
	frame=single,
	framesep=2mm,
	breaklines=true,
	breakanywhere=true
}
% 使章节标题左对齐（同时保留 ctex 默认样式）
\ctexset{section={format+=\raggedright}}
%1.3倍行距
\renewcommand{\baselinestretch}{1.3}
% 数学左对齐缩进为0
\setlength{\mathindent}{0pt}
% 增大矩阵的行距
\renewcommand{\arraystretch}{1.3}
\pagestyle{empty}
\begin{document}
\section*{题目1}

\noindent \textbf{信号量初始值：}

$a = 2, \quad b = 0$

\vspace{1em}

\textbf{t1 的代码：}
\begin{itemize}
    \item [{[1]}] \verb|P(a); P(a);| 
    \item [{[2]}] \verb|V(b); V(b);| 
\end{itemize}

\textbf{t2 的代码：}
\begin{itemize}
    \item [{[1]}] \verb|P(b);| 
    \item [{[2]}] \verb|V(a);| 
\end{itemize}

\section*{题目2}

\textbf{doit 函数实现：（\&a[M*i]这个地址值不会随着i的改变而改变）}
\begin{verbatim}
void doit(void *args) {
    int *arr = (int *)args;
    for (int i = 0; i < M; i++) {
        Hash_Insert(&hashtb, arr[i]);
    }
}
\end{verbatim}

\end{document}